\chapter{随机森林与 Bagging}
\chaptergoal{解释为何组合多个不完全相同的模型可能改善预测；区分 Bootstrap、Bagging 与随机森林的随机特征选择；使用 sklearn 比较单树与森林；理解袋外评价和特征重要性的适用范围。}

前一章的决策树容易解释，但单棵树对训练样本和局部划分比较敏感。若在训练数据略有变化时，树的顶部分裂条件就发生改变，预测可能相差明显。随机森林不是把一棵树训练得越来越深，而是训练多个有差异的树，再综合它们的预测。这里的关键是理解\textbf{差异从哪里来、集成怎样工作，以及这种差异是否有利于预测}。
\section{单棵树的变化与集成思想}

\textbf{集成学习（ensemble learning）}是把多个模型的预测按照明确规则组合成一个总体预测的方法。被组合的每个模型称为\textbf{基学习器}或基模型。“集成”不是把训练数据的不同列随意合并，也不是挑其中一棵分数最高的树；它最终使用多个模型共同作出预测。

假设三棵回归树对同一条新记录分别预测10、12、14，等权平均得到12。如果三棵树每次都给出完全一样的错误，平均不会消除错误；如果它们的偶然偏差有所不同，平均可能减少预测的波动。这里所说的\textbf{预测方差}是指在训练数据等条件改变时预测值的波动程度，与目标数据自身的方差不是同一对象。

考虑一批固定的 Wine 记录，随机抽取不同训练子集，每次训练一棵同样配置的决策树。某些分裂条件可能反复出现，另一些则对抽样变化非常敏感。如果只采用一棵树，预测依赖这一次训练样本的偶然性。Bagging 则从多次随机重采样开始，训练多棵树，最后取多数投票或平均值。

对于回归，设第 $b$ 棵树的预测是 $f_b(\mathbf x)$，Bagging 的平均预测是
\begin{equation}
\hat f(\mathbf x)=\frac1B\sum_{b=1}^B f_b(\mathbf x).
\end{equation}
如果各棵树的误差不是完全同向变化，平均可能减小预测波动。这里的“减小”并非对一切数据都保证，树之间的相关程度、训练样本数量、偏差与噪声都会影响收益。分类采用投票或者平均类别概率等组合规则，具体结果也需核对实现方式。

\subsection{从平均预测推导波动的变化}
为了把“多个模型取平均”说得更准确，固定一个待预测输入 $\mathbf x$，把在不同训练抽样或算法随机状态下得到的第 $b$ 个预测记为随机量 $Z_b$。本节讨论这些预测的波动，而不是把不同测试对象的预测值混在一起计算方差。一般有
\[
\operatorname{Var}\!\left(\frac1B\sum_bZ_b\right)
=\frac1{B^2}\left[\sum_b\operatorname{Var}(Z_b)+2\sum_{b<c}\operatorname{Cov}(Z_b,Z_c)\right].
\]
若为了理解机制，作简化假设：各模型方差均为 $\sigma^2$，任意两个不同模型的相关系数均为 $\rho$，则
\begin{equation}
\operatorname{Var}(\bar Z)=\sigma^2\left[\rho+\frac{1-\rho}{B}\right].
\label{eq:forest-variance}
\end{equation}
推导只需数清楚：有 $B$ 个自身方差项，和 $B(B-1)$ 个有序的不同模型协方差项。先除以 $B^2$，再整理即可。实际森林不满足每对相关程度完全相同，式子在这里是有明确假设的解释模型，不是用于直接预测每个实验分数的定律。

\worked{加树数量与减小相关性，并不是同一个操作}
取 $\sigma^2=1$。若 $\rho=0$，10棵树平均后的方差为0.1，100棵为0.01。若 $\rho=0.2$，相应值为0.28和0.208；树再多，公式中的相关部分0.2仍然存在。若 $\rho=0.8$，100棵平均后的方差为0.802，改善很有限。

这个例子说明：复制同一棵树一百次并平均，不会获得一百份独立信息；数量的收益取决于误差是否共同变化。随机特征选择就是试图减少某些共同的划分偏好，但减少相关性同时可能使单棵树变弱，因此不能把“越随机越好”当作结论。

\fig{v3_ch09_correlation.pdf}{在相同方差、相同两两相关性的人工模型下，预测平均的理论方差与模拟检查。图只检验平均的统计机制，不是葡萄酒数据的实际误差曲线。}

还需要区分\textbf{方差}与\textbf{偏差}。若每棵模型平均都偏高2个单位，等权平均仍偏高2个单位。集成可以减少部分偶然波动，却不会自动修正所有模型共同缺少的物理变量或错误的标签。第4章平方误差的偏差—方差分解，在这里给出了理解平均收益的框架。

\subsection{分类投票为什么也需要差异}
如果五个分类器始终在同一批样本上犯同样的错，多数投票不会纠正这些错误。相反，若某个样本上三个分类器正确、两个错误，投票可以优于其中某一个错误分类器。但是这种逐样本现象不等于“单模型准确率超过一半就必然靠投票变好”：模型之间的错误依赖、类别数、概率聚合方式和样本分布都影响最终结果。

本书把分类组合的直觉与回归平均的方差公式分开讲述。不能将式\eqref{eq:forest-variance}中的方差下降，直接解释为分类准确率增加同样百分比。分类还涉及最后从概率或票数变为类别的非线性决策步骤。

\section{有放回重采样与 Bagging}

\textbf{有放回抽样}指抽取一条记录后，把它放回可抽取集合，下次仍可能再次抽到。\textbf{自助采样（Bootstrap）}在这里用这一方式，从已有训练集生成与它同样大小的重采样集。\textbf{Bagging}是Bootstrap aggregating的简称：分别在多个自助样本上训练模型，再聚合预测。

原训练集若只有A、B、C、D、E五条记录，一次抽五次可能得到 $[A,C,C,E,A]$。A与C各出现两次，B与D没有出现。该树的袋外记录就是B和D；另一棵树可抽得 $[B,E,D,D,C]$，它的袋外记录为A。两棵树的训练数据既有重叠，也有各自缺席的记录。因此“不同重采样”不等于“完全不重叠的真实实验”。

Bootstrap 从 $n$ 条训练记录中有放回地抽取 $n$ 次，形成一个训练子集。一条原始记录可能被抽到多次，也可能一次未被抽到；后者称为该树的袋外记录。不同重采样训练集提供了模型之间的差异。即使每棵树用相同的算法和参数，训练记录的差别也能使预测不同。

\textbf{Bagging} 是一种组合策略，不限于决策树；随机森林则以树为基本学习器，并进一步在节点划分时引入随机特征选择。Bagging 并不必然把训练偏差降到很小，它的主要价值之一是降低某些高波动模型的方差。对本身非常稳定的基学习器，Bagging 收益可能不明显。

\subsection{为什么一份自助样本没有包含所有原记录}
从 $n$ 条记录中有放回地抽取 $n$ 次。固定一条原记录，每次没抽到它的概率为 $1-1/n$，所有抽取都没抽到的概率为
\[
P(\text{未出现})=(1-1/n)^n.
\]
因此原记录至少出现一次的概率为 $1-(1-1/n)^n$，自助样本中\textbf{不同原记录数量的期望}为 $n[1-(1-1/n)^n]$。当 $n$ 较大时，未出现比例接近 $e^{-1}\approx0.368$，出现比例接近0.632。这里说的是期望，不是每棵树都精确使用63.2\%不同记录。

\worked{五条记录抽五次}
对 $n=5$，一条记录未出现概率为 $(4/5)^5=0.32768$；不同记录数量的期望为 $5(1-0.32768)=3.3616$。虽然每份自助样本都含5次抽取，独特记录平均只有约3.36条。重复出现的记录相当于在该树训练中获得较高的重复权重。若每次不放回抽取全部5条，就只是改变顺序，不会获得这种样本组成的差异。

这也解释了为什么不能把Bootstrap称为“生成新的独立观测”。它重新组织已有数据，没有增加测量信息。它可以近似研究在当前经验分布下的抽样变化，也可制造基模型差异；能够支持什么结论，仍受原样本代表性和独立性限制。

\subsection{把Bagging的训练与预测写成算法}
训练阶段先固定基学习器配置和树数 $B$。对 $b=1,\ldots,B$，独立产生一组重采样索引，在相应输入与目标上训练一个模型，并保存索引与模型。预测阶段，对同一个新输入分别调用所有已训练模型，再按规定聚合。这里的“独立产生索引”并不意味着所有树的预测统计上独立，因为它们来自同一原始数据集。

\begin{lstlisting}[caption={显式展示三棵自助树，便于核对重采样与预测}]
import numpy as np
from sklearn.tree import DecisionTreeRegressor
rng = np.random.default_rng(9)
models, sampled_rows = [], []
for b in range(3):
    idx = rng.integers(0, len(X_tr), size=len(X_tr))
    m = DecisionTreeRegressor(max_depth=3, random_state=b)
    m.fit(np.asarray(X_tr)[idx], np.asarray(y_tr)[idx])
    models.append(m)
    sampled_rows.append(idx)
predictions = np.column_stack([m.predict(X_val) for m in models])
mean_prediction = predictions.mean(axis=1)
\end{lstlisting}

这个代码用于已准备好的回归训练与验证数据。不能把类别代码送入它再平均代码数值来代替分类；分类需要明确投票或概率平均。本节演示显式循环，正式实验可使用 \texttt{BaggingClassifier} 或 \texttt{BaggingRegressor}，其内部还处理并行、随机状态和其他配置。

\section{随机特征选择与随机森林}

\textbf{随机森林（random forest）}将许多决策树组合起来，并在训练树时引入随机性，使各树不要总作相同判断。在本章的典型配置中，一种随机性来自样本重采样，另一种来自每个节点尝试划分时的候选特征选择。例如有13个输入时，一个节点可能仅比较其中3个随机选出的特征；下一个节点可以重新选择候选，而不是一棵树始终只使用同3列。

普通 Bagging 中各棵树虽然来自不同重采样，但如果某个特征非常强，它们的顶部分裂仍可能高度一致。随机森林在每个节点寻找划分时只考虑部分随机选择的特征，从而增加树之间的差异。实际实现中 \texttt{max\_features} 控制节点候选特征数，\texttt{bootstrap} 控制是否采用有放回采样。

\texttt{RandomForestClassifier} 和 \texttt{RandomForestRegressor} 分别提供分类与回归。分类森林的预测不是把每棵树的类别数字做普通平均；实现通常综合各树的类别概率再作出分类决定。

以三棵树为例，它们给类别A的概率分别为0.51、0.51、0.01。按每棵树的最终类别投票，A得到两票；按概率平均，A的平均概率却只有 $1.03/3\approx0.3433$，因此选择B。两种聚合规则并不等价。sklearn的分类随机森林采用平均类别概率，再取最大者；本书解释程序输出时依据这一规则。\textsuperscript{[8]}
回归森林则对连续预测求平均。因此，读懂森林输出需要同时说明预测目标的类型和聚合规则。

\subsection{从强特征主导到节点级随机候选}
假设有一个测量指标在多数自助样本上都略优于其他指标。普通Bagging的许多树可能都以它作根节点，随后沿相似方向细分。如果该指标含有某种系统误差，多棵树的错误也可能一致。随机森林在节点比较中限制候选特征，有时迫使树利用其他有预测作用的信息，从而改变模型之间的相关性。

具体地，若13个特征中每个节点随机选择3个候选，并不表示最终树只能使用3个特征。一条较长路径或整棵树可以在不同节点用到更多列。也不表示对选出的3列取随机阈值；普通随机森林仍在允许特征上搜索较好的划分。另一些极端随机树方法进一步随机化阈值，其机制不能与这里的配置混为一谈。

\texttt{max\_features}较小时，单个节点可比较的信息减少，树间差异可能增大；但若真正有用的特征非常少，节点频繁错过它们，单树偏差也可能增加。较大的候选集合提升单次划分能力，却可能让树更相似。应通过固定其他条件的实验观察这种权衡，而不是把默认值当作科学结论。

\subsection{叶频率到森林概率的完整过程}
对类别 $k$，第 $b$ 棵树把样本送入某叶，给出叶中的类别频率 $p_{bk}(\mathbf x)$。分类森林常使用
\[
\bar p_k(\mathbf x)=\frac1B\sum_b p_{bk}(\mathbf x),\qquad
\hat y=\arg\max_k\bar p_k.
\]
这些概率先逐树求出，再对每类平均；不是先把所有树的叶样本记录合成一张大表重新统计，也不是平均0、1、2这样的类别编号。某些树的叶只有少量训练记录，另一些叶较大，sklearn的普通概率聚合并不会因此按叶样本数重新加权每棵树。

模型给出0.9并不自动意味着所有这类预测中恰有90\%正确。概率\textbf{校准}是另一项需要利用未参与拟合的数据检查的性质；增加树数可能使概率更平滑，却不保证校准。第6章已经区分概率和决策，第16章会把概率可靠性与错误分析联系起来。

\section{树的数量和单树复杂度}
\texttt{n\_estimators} 控制树数量。树太少时，集成结果可能随随机种子产生较大波动；增加树数常使表现趋于稳定，但训练及预测成本也会增长。树数增加不保证验证误差持续下降，更不能补救错误的目标定义或数据泄漏。

\texttt{max\_depth}、\texttt{min\_samples\_leaf} 等控制每棵树的复杂度。浅树可能无法捕捉较复杂关系，深树可能对数据噪声敏感。森林通过集成减小部分不稳定性，但仍需通过验证调整复杂度。教学实验应分别改变“树的数量”和“单树复杂度”，避免同时改变多个参数后无法判断原因。

\begin{lstlisting}[caption={Wine：浅树与随机森林比较}]
from sklearn.datasets import load_wine
from sklearn.model_selection import StratifiedKFold, cross_val_score
from sklearn.tree import DecisionTreeClassifier
from sklearn.ensemble import RandomForestClassifier

X, y = load_wine(return_X_y=True)
from sklearn.model_selection import train_test_split
X, X_hold, y, y_hold = train_test_split(
    X, y, test_size=0.2, stratify=y, random_state=42
)
# All comparisons below use development X, y only.
cv = StratifiedKFold(5, shuffle=True, random_state=7)
models = {
    "tree": DecisionTreeClassifier(max_depth=4, random_state=42),
    "forest": RandomForestClassifier(n_estimators=150,
                                      min_samples_leaf=2,
                                      random_state=42, n_jobs=1),
}
for name, model in models.items():
    s = cross_val_score(model, X, y, cv=cv)
    print(name, s.mean(), s.std())
\end{lstlisting}



\subsection{一项只改变树数的重复实验}
配套实验在同一葡萄酒开发划分上固定叶最小样本数与其他参数，只改变树数为10、30、100、300，每个设置使用5个预先列出的随机种子。结果如下；均值和标准差来自随机种子重复，不是交叉验证折间波动。
\begin{center}\small
\begin{tabular}{rrr}\toprule
树数&36条验证记录上的准确率均值&种子间样本标准差\\\midrule
10&0.9778&0.0232\\
30&0.9889&0.0152\\
100&1.0000&0.0000\\
300&1.0000&0.0000\\\bottomrule
\end{tabular}\end{center}
\fig{v3_ch09_tree_count.pdf}{固定训练与验证记录时，不同树数下5个随机种子的结果。准确率达到平台，不说明概率完全相同，也不代表对更大总体没有误差。}

100与300棵都把这36条记录分对，因而本指标无法在两者之间提供进一步分辨。不能据此声称“300棵严格更好”，也不能把标准差0当作预测不确定性0。这里固定了训练样本和验证样本，改变的只是算法随机性；换一批真实训练记录可能产生更大变化。

\subsection{控制模型规模时应记录的量}
树数增加主要降低有限集成的随机波动，并增加存储和预测成本；单树深度或叶样本数改变的是每棵树的拟合方式。两者不能相互替代。1000棵过浅的树可能仍有较大共同偏差，10棵极深树也可能不稳定。除准确率外，可记录实际总节点数、训练时间和一次批量预测时间；在比较硬件前必须固定数据、线程数和测量方法。

随机森林的并行训练可以提高速度，但多线程不等于改变统计目标。实验为了容易重复默认 \texttt{n\_jobs=1}；在自己的机器上提高线程数时，应把资源设置记入记录。不要把不同模型默认线程数造成的差异直接解释成算法本身的速度优势。

\section{袋外评价的基本含义}

对一条记录做袋外预测时，只允许那些训练重采样中没有抽到它的树参加平均或投票。例如A未进入第2、第5、第7棵树，就只汇总这几棵树对A的预测，再与A的已知目标比较。不能先让所有树预测A，再把这个结果称为袋外评价，因为其中一些树已经用A训练过。

一棵使用 Bootstrap 训练的树，没有见过部分原始记录；这些袋外记录可用于衡量该树在非训练记录上的表现。汇总每条训练记录在相应袋外树上的预测，就可构造森林的袋外评价（out-of-bag, OOB）。这减少了另行划出验证集的一部分需要，但 OOB 与交叉验证具有不同的数据使用结构；它也不能代替最终独立测试。

袋外记录的身份来自重采样，并不自动尊重科学对象分组。如果同一天体的另一条观测被抽进训练，即便当前观测对某棵树是袋外，仍可能发生对象信息泄漏。因此 OOB 的“袋外”不是“在科学意义上完全独立”。

\subsection{逐样本组成袋外预测}
设第 $i$ 条记录未出现在树集合 $\mathcal O_i$ 的训练自助样本中，则回归袋外预测为
\[
\hat y_i^{\mathrm{OOB}}=\frac1{|\mathcal O_i|}\sum_{b\in\mathcal O_i}f_b(\mathbf x_i).
\]
分类对应平均相关树的类别概率。每条记录使用的树集合不同、数量也可能不同。若树数太少，某些记录可能没有任何袋外树，就无法形成这样的评价；应检查实现警告和有效覆盖，而不是只读取一个返回分数。

\worked{一条记录不能让见过它的树参与袋外核对}
假设四棵树对记录A的类别1概率分别为0.9、0.2、0.8、0.4，其中只有第2和第4棵没有训练过A。袋外概率是 $(0.2+0.4)/2=0.3$，不是四棵平均0.575。若真实类别为1，前者在0.5阈值下是一次袋外错误；后者可能因为混入记住A的树而看起来正确。这一简单计算说明，“袋外”是对预测来源的约束，不是训练后任意贴的标签。

当前葡萄酒106条训练记录、300棵树的实验中，袋外准确率约0.9623，另36条固定验证记录的准确率为1.0000。一条训练记录对单棵树未入袋的理论概率为 $(1-1/106)^{106}\approx0.3661$。袋外和验证分数不同并不矛盾：它们使用不同记录及不同预测集合，也具有有限样本波动。

\subsection{袋外评价与完整流水线的边界}
若先在全部开发数据上用目标挑选特征，再训练森林并报告OOB，特征选择已经见过那些所谓袋外记录的标签。森林内部的重采样不能撤销外部发生的信息泄漏。对需要学习的预处理、目标驱动特征选择或复杂调参，使用整体交叉验证往往更容易正确表示完整流程。

同一对象的多条记录也是类似问题：当前行没入袋，不代表同一对象的其他行没入袋。若目标是推广到新对象，应以分组验证作为主要评价。OOB是一种有用的内部诊断，不是无需理解数据结构的独立性保证。

\section{葡萄酒分类与住房回归的模型比较}
分类用 Wine 比较单树与森林，记录五折均值和波动；回归用住房数据比较常数基线、回归树和随机森林，重点观察 MAE/RMSE 及不同目标区间的残差。树模型一般不需要为阈值划分进行相同形式的标准化，但为了公平比较，必须保持可用特征集合一致。

\begin{lstlisting}[caption={回归森林的可重复配置，数据由本书第5章准备}]
from sklearn.ensemble import RandomForestRegressor
reg = RandomForestRegressor(n_estimators=200,
                            min_samples_leaf=3,
                            random_state=42, n_jobs=1)
# reg.fit(X_dev, y_dev)
# Evaluate on development folds before one final holdout test.
\end{lstlisting}

图示中的运行时间与参数效果取决于软件和硬件版本，不能未经实测就写入“森林训练必然更快/更慢”的结论。

\subsection{固定五折后的实际比较}
为了单独观察组合方式，本轮在142条葡萄酒开发记录上使用相同5折。分别拟合不限制深度的单树、100棵自助树的Bagging，以及100棵采用节点随机候选特征的森林。模型配置由实验文件明确保存，最终36条保留数据未用于本表。
\begin{center}\small
\begin{tabular}{lrr}\toprule
模型&五折准确率均值&折间标准差\\\midrule
单棵较深树&0.9227&0.0261\\
Bagging&0.9510&0.0278\\
随机森林&0.9860&0.0172\\\bottomrule
\end{tabular}\end{center}
本次结果与“不同树的组合有帮助”的预期一致，但三个数不能证明任何数据上的固定排名。特别是Bagging折间标准差略高于单树，不能因为均值提高就声称所有稳定性指标都改善。正确解释要区分平均表现、划分间波动、算法种子波动以及误差发生在哪些对象上。

如果要报告森林相对单树的改善，优先在同一折内相减，检查差异是否在多个折都一致，而不是把来自不同划分的最高分相减。小样本五折不是五个独立科研实验；无需为了得到一个“显著”结论而忽略训练数据重叠。

\subsection{回归的平均也不自动提供可靠区间}
回归森林返回树预测的平均。可以计算各树预测的离散程度作为模型内部诊断，但这些树相关，而且通常不包含未来测量噪声的完整贡献。直接把“树间标准差乘2”当作95\%预测区间，没有一般保证。专业任务需要区分模型平均预测的变化与下一次观测值的散布，第16章和附录E再讨论这一区别。

住房和风洞数据的实战应继续使用第5章的目标、单位、开发划分和基线。在无法获取外部数据的离线环境，书中共同的人工响应实验仍能比较代码和机制；报告必须写“人工响应”，不能把它的RMSE改名成住房价格误差。

\section{特征作用与专业迁移}
树模型提供基于不纯度或误差减少的特征重要性，但这种重要性可能偏向取值较多、可切分方式丰富的特征，也可能在相关特征之间分配得不稳定。若某一测光颜色最常出现于分裂条件，不代表该颜色是天体类别的单一物理原因。第16章将进一步使用置换重要性，并讨论相关特征与分组误差。

\fig{ch09_rf_importance.pdf}{Wine 分类森林的示例特征重要性。该重要性依赖当前训练数据、模型与重要性定义，不直接表示因果效应。}

在专业迁移中，可复用 SDSS 或 Sonar 分类、Airfoil 回归，让森林与第6章逻辑回归、第8章决策树在一致的开发验证方案下比较。每个实验至少分析一种具体错误或残差结构，而不只报告单一准确率。

\subsection{不纯度重要性究竟统计了什么}
对某个特征，把整棵树中由它引起的节点加权不纯度下降相加，再在所有树之间平均并归一化，就得到常见的不纯度重要性。它描述该特征在当前训练搜索中承担了多少拟合改善，不是直接度量将该物理量人为改变后目标会变化多少。

若一个随机编号列几乎每行都不同，它有很多可比较阈值，有机会偶然获得训练改善；只有两个取值的真实特征候选机会较少。因此重要性可能受到可选划分数量影响。若两个特征高度相关，算法可以在两者之间选择，重要性可能被分摊，也可能因为某次局部选择集中在其中一列。低重要性不必然意味着该观测量没有科学信息。

应先确认模型在验证数据上确有预测能力，再解释其重要性。第16章的置换检验会问“破坏某列与对象的对应后，验证表现下降多少”，这是另一个问题。两种重要性可以共同分析，但不能要求数值或排序必须一致。

\subsection{一项可完成的专业比较}
对声呐回波分类，先保留已定义的60项输入和对象范围，比较标准化逻辑回归、浅树与森林。树无需为距离而缩放，不表示逻辑回归也应取消缩放；公平性要求信息一致和验证合规，不要求每个模型内部步骤完全一样。对风洞声压级预测，增加按工况留出的比较，查看森林在新工况上的误差是否与随机行划分不同。

最终解释可以是：“在这些开发折上，森林平均比单树少错若干条，主要改善某两类对象的混淆；重复对象已按组处理，跨仪器仍未检验。”这类具体结论比“森林更智能，因此发现了物理规律”更有用。算法复杂程度是完成任务的手段，不是科学结论强弱的尺度。

\chapterpractice{分别改变森林的树数和叶节点规模，并与单棵树比较；在 Wine 中观察训练子样本变化导致的结果波动；在住房或专业回归数据上比较森林与线性模型的残差。}
\chapterexercise
\begin{enumerate}
\item 为什么 Bootstrap 子集可能包含重复记录？每棵树是否使用完全不同的样本？
\item Bagging 与随机森林的区别是什么？其随机性来自哪些环节？
\item 袋外记录能否保证同一天体的不同观测之间没有信息泄漏？
\item 为什么某个特征的重要性较高，并不足以证明该特征的变化导致目标发生变化？
\end{enumerate}

\section*{深化练习与方法检验}
本组题强调解释、计算与验证；D编号在本章内独立计数。计算题应写出所用定义和中间步骤，实验题保留实际输出，包括不支持预期的结果。
\begin{enumerate}[label=D\arabic*.]
\item 假设 $B$ 个模型的预测误差方差均为 $\sigma^2$，两两相关系数均为 $\rho$。推导等权平均误差的方差，并以 $B=10,\rho=0.2,\sigma^2=1$ 计算数值。
\item 从 $n$ 条记录有放回抽取 $n$ 次，求某条记录未被抽中的概率，并解释极限 $e^{-1}$ 的含义。它能否自动证明袋外样本彼此独立？
\item 三棵树对类别A的概率为0.51、0.51、0.01。比较硬投票和平均概率的分类；为什么不能把类别数字直接当连续预测平均？
\item 森林对一个特征的重要性很高。设计一项检验以区分“该特征独有信息”和“与其他特征可替代的信息”，并说明它能否证明物理因果。
\end{enumerate}
